\section{Coefficient fields of CM modular forms}
\label{sec:cm}
A normalized newform $f=\sum_{n\geq1}a_nq^n$ has a coefficient field
$K_f=\mathbb Q(a_1,a_2,\ldots)$. Fixing its weight and the degree of this
field still allows the level to grow without bound. Martin's Conjecture
1(a) asks whether only finitely many fields occur at fixed weight and
degree for trivial nebentypus (the Dirichlet character in the modular
transformation law). A family on which this can be decided consists of
the forms with complex multiplication, abbreviated CM. The decisive
zeros belong to quadratic Dirichlet $L$-functions. One needs no
hypothesis on the zeros of the modular forms' own $L$-functions.

\paragraph{Finiteness with varying character.}
Take the repository's Dirichlet GRH theorem as input. For fixed integers
$k\geq2$ and $d\geq1$, only finitely many isomorphism classes of fields
$K_f$ occur among CM newforms of weight $k$ with
$[K_f:\mathbb Q]=d$, even as both the level and nebentypus vary.

Martin proves the trivial-character case under quadratic Dirichlet GRH
(Theorem~1). His Remark~1.1 anticipates a fixed-character analogue and
points to the same methods for CM forms. We give the simultaneous
deduction, including the field identity that allows the character to
vary. The class-group and value-field arguments are Martin's.
\href{https://arxiv.org/pdf/2509.24119v2}{Martin, \emph{Rationality fields of CM modular forms}, Remark 1.1 and Sections 2--3}.

\paragraph{Recover the Hecke character's field.}
Write $E$ for the imaginary quadratic CM field and $\ell=k-1>0$.
The CM correspondence realizes $f$ as the theta series of an algebraic
Hecke character $\psi$ of $E$ of type $(\ell,0)$ and conductor
$\mathfrak m$. Thus $a_n$ is the sum of $\psi(\mathfrak a)$ over integral
ideals of norm $n$, extending $\psi$ by zero at ideals meeting the
conductor. Let $L$ be the field generated by its values and $K=K_f$.
Then $K\subset L$. We claim that $L=EK$ for every nebentypus.

For $p=\mathfrak p\bar{\mathfrak p}$ split and prime to the conductor,
put $u=\psi(\mathfrak p)$ and $v=\psi(\bar{\mathfrak p})$. The
coefficients at $p$ and $p^2$ give
\[
       u+v=a_p,\qquad uv=a_p^2-a_{p^2}.
\]
Take $h$ to be the exponent of the ray class group modulo $\mathfrak m$.
Then $\mathfrak p^h=(\alpha)$ and
$\bar{\mathfrak p}^{h}=(\beta)$ with generators congruent to $1$
modulo $\mathfrak m$. Hence $u^h=\alpha^\ell$ and $v^h=\beta^\ell$
belong to $E$. Their principal ideals are distinct, and neither is
invariant under conjugation. Consequently $u^h$ generates $E$ over
$\mathbb Q$, so $E\subset L$. Every automorphism of
$\overline{\mathbb Q}$ fixing $EK$ preserves the pair of roots $u,v$ of
$X^2-a_p X+(a_p^2-a_{p^2})$ and fixes their distinct $h$th powers.
It therefore fixes each root. This gives $u,v\in EK$. At a split prime
where only one factor is prime to $\mathfrak m$, its value equals $a_p$.
At an inert prime the sole prime ideal has value $a_{p^2}$, and at a
ramified prime its value is $a_p$, whenever it is prime to the conductor.
Prime ideals generate the ideal group, which proves the claim. In
particular $[L:E]=d/[K\cap E:\mathbb Q]$ divides $d$.

\paragraph{Bound the CM discriminant.}
For an ideal $\mathfrak a$ prime to $\mathfrak m$, write
$\mathfrak a^h=(\gamma)$ with $\gamma\equiv1\pmod{\mathfrak m}$.
Taking $h$th powers of ideals and using $\psi(\mathfrak a)^h=\gamma^\ell$
gives
$(\psi(\mathfrak a))=(\mathfrak a\mathcal O_L)^\ell$ in $L$.
The ideal norm from $L$ to $E$ shows that
$\mathfrak a^{\ell[L:E]}$ is principal. Since every class has such a
representative, $e_E=\exp\mathrm{Cl}(E)$ divides $(k-1)d$.

Here is where the analytic result enters. Write $D$ for the fundamental
discriminant of $E$. If $p=\mathfrak p\bar{\mathfrak p}$ splits, take
$\mathfrak p^{e_E}=(\alpha)$. Unequal valuations at the two primes above
$p$ force $\alpha\notin\mathbb Q$. Writing
$\alpha=(a+b\sqrt D)/2$ with $b\ne0$ gives
\[
       p^{e_E}=|N_{E/\mathbb Q}\alpha|\geq |D|/4 .
\]
A sufficiently small split prime therefore bounds the discriminant of every possible $E$.
A numerical version gives a test for an individual proposed CM field. Put
$T=\log|D|$ and $M=(k-1)d$. If $T>25$, Theorem~3 of Elsenhans,
Kl\"uners and Nicolae gives $p\leq(1.881T+6.18)^2$ under ERH for the
quadratic field. Combining it with the preceding ideal-norm inequality
forces
\[
               T\leq\log 4+2M\log(1.881T+6.18).
\]
The right side grows only logarithmically with $T$, so this is an
explicit discriminant cutoff for every weight and coefficient degree.
For example, in weight two and degree seven the class-group exponent
divides seven, and their Table~2 yields $|D|\leq 3.9\cdot10^{30}$.
\href{https://arxiv.org/pdf/1803.02056}{Elsenhans, Kl\"uners and Nicolae, Theorem~3 and Table~2}.

\paragraph{Finish the field count.}
For any one of these $E$, Martin's Corollary~3.6 gives finitely many
value fields $L/E$ of each bounded degree, uniformly in the conductor.
Indeed the image on principal ideals generates $E(\mu_r)$, where $r$
is the order of the finite modulus character. Thus
$\varphi(r)\leq [L:\mathbb Q]\leq 2d$. Choose ideal-class generators
$\mathfrak t_i$ with $\mathfrak t_i^{n_i}=(\theta_i)$.
Martin's Proposition~3.4 places $L$ in the field obtained by adjoining
$\mu_{r e_E}$ and the $n_i$th roots of $\theta_i^\ell$ to $E$, independently
of $\mathfrak m$. There are finitely many possible $r$, and each resulting
number field has only finitely many subfields. This proves the assertion.

\paragraph{Why weight two admits more.}
For trivial nebentypus Martin's Theorem~2 lists 23 real quadratic and 18 cubic
fields using quadratic Dirichlet GRH. In weight two the divisibility
$h_E\mid[L:E]\mid d$ of Yang and Gonz\'alez gives an independent, unconditional
reduction from class-group exponent to class number. Gonz\'alez's Theorem~3.1
already proves weight-two finiteness in all degrees and for all characters.
Our interest here is the exact list over all levels. The next sections give the
genus-field and cubic arguments, then the classification in higher prime degrees.
For general weights the split-prime proof above works whenever a
sufficiently small split prime is available. The full Dirichlet GRH supplies
one route. Martin's conjecture for non-CM newforms requires other ideas.
\href{https://arxiv.org/pdf/2509.24119v2}{Martin, Theorems 1--2};
\href{https://ems.press/content/serial-article-files/38349}{Gonz\'alez, Theorems 3.1--3.2}.
